\section*{第 \ref{ch:g5:brain-in-command} 章——大脑指挥全身}

\begin{solution}{exo:g5:brain-in-command:1}
脑：接收全部汇报，做出全部决定，发出全部指令。
\omterm{def:g5:brain-in-command:system}{脊髓}：从脑向下，穿过\omterm{def:g4:classifying-animals:vertebrate}{脊柱}内部的一条粗缆线。
\omterm{def:g5:brain-in-command:system}{神经}：呈\omterm{def:g1:plants-around-us:tree}{分支}状的细丝，在脑或\omterm{def:g5:brain-in-command:system}{脊髓}与
每个器官之间传递信息。
\end{solution}

\begin{solution}{exo:g5:brain-in-command:2}
\omterm{def:g1:the-five-senses:sense}{感觉器官}向脑汇报；脑作出决定并发出指令；
\omterm{def:g5:brain-in-command:system}{神经}把指令传给\omterm{def:g3:bones-and-muscles:muscle}{肌肉}，使\omterm{def:g3:bones-and-muscles:muscle}{肌肉}收缩。
\omterm{def:g1:the-five-senses:sense}{感觉器官}汇报，脑做决定，
\omterm{def:g3:bones-and-muscles:muscle}{肌肉}服从指令。
\end{solution}

\begin{solution}{exo:g5:brain-in-command:3}
脑位于\omterm{def:g3:bones-and-muscles:skeleton}{颅骨}内；
\omterm{def:g5:brain-in-command:system}{脊髓}位于由一块块小骨叠成的\omterm{def:g4:classifying-animals:vertebrate}{脊柱}内。
\end{solution}

\begin{solution}{exo:g5:brain-in-command:4}
信号到身体作出反应之间的延迟，对熟练的人完成简单任务，
约为五分之一秒，
儿童则稍长一些。
\end{solution}

\begin{solution}{exo:g5:brain-in-command:5}
足部\omterm{def:g1:the-five-senses:sense}{皮肤}汇报（触觉和疼痛），\omterm{def:g5:brain-in-command:system}{神经
系统}决定缩回，腿部\omterm{def:g3:bones-and-muscles:muscle}{肌肉}服从指令，
猛地把脚抽回。（如此快速的缩回，部分由\omterm{def:g5:brain-in-command:system}{脊髓}
自身处理——相当于指挥中心的
应急台。）
\end{solution}

\begin{solution}{exo:g5:brain-in-command:6}
它昼夜不停地维持\omterm{def:g4:breathing:movements}{呼吸}，并根据身体的需要调节
心跳的节奏（此外还负责：储存记忆、
做梦）。
\end{solution}

\begin{solution}{exo:g5:brain-in-command:7}
\omterm{prop:g1:healthy-habits:needs}{睡眠}——这是整理与修复的时间——也占用
身体很大且稳定的一部分\omterm{prop:g3:balanced-diet:jobs}{能量}和氧气。二者缺一，它
反应变慢，注意力分散，判断也变差。
\end{solution}

\begin{solution}{exo:g5:brain-in-command:8}
答案不唯一。通常的接住距离约为 $10$--$20$ 厘米；
多试几次，数值通常会略有改善并趋于稳定。
\end{solution}

\begin{solution}{exo:g5:brain-in-command:9}
因为每次接住都必须经过三个步骤——眼睛的
信息要传送，脑要计算并作出决定，
指令要传回，\omterm{def:g3:bones-and-muscles:muscle}{肌肉}要收缩——每一步
都需要真实的时间。熟练可以缩短这一循环，
却无法将其去掉。
\end{solution}

\begin{solution}{exo:g5:brain-in-command:10}
因为最好的刹车只能在驾驶员的这一循环
完成之后才起作用——高速行驶时，车辆在这段时间里约前进 \qty{7}{m}，
指令“刹车！”才传到脚上。两车之间的距离，
正是为这一循环留出的空间。注意力可以缩短汇报与决定的步骤；
分心的\omterm{def:g5:brain-in-command:system}{大脑}则会多耗掉几米。
\end{solution}

\begin{solution}{exo:g5:brain-in-command:11}
\omterm{def:g5:brain-in-command:system}{大脑}的任务——决定与发令：练习
铺设了更好的通路，于是同样的信息
能快速、轻松地被送往正确的指令。\omterm{def:g3:bones-and-muscles:muscle}{肌肉}训练
增强的是\omterm{def:g3:bones-and-muscles:muscle}{肌肉}力量；技能训练增强的则是\omterm{def:g5:brain-in-command:system}{大脑}的信息分配。
一个锻造动力，另一个锻造指挥。
\end{solution}
